% ch11 Conclusions (书页282-283 / p300-301)
\chapter{结论}

全书的论证在最后一章收拢。原书结论没有新公式、新推导，而是把从股票期权到远期利率的全部内容作一次回顾与定位；本导读据此归纳为三层意思。

第一层：量子理论的两种表述各司其职。大多数数理金融模型只涉及有限个自由度（单一标的、即期利率等），它们用随机微积分能处理，但用路径积分与哈密顿量同样能处理，而且路径积分完全有限、可以算出所有感兴趣的量——股票期权的定价核（pricing kernel，即条件概率）就是最直接的例子。哈密顿表述则带来了随机微积分之外的新工具：势（potential）与本征函数（eigenfunction）使障碍期权（barrier option）与路径依赖（path-dependent）期权的建模有了统一的几何图像，完备性方程（completeness equation）给出计算定价核的系统办法，而“鞅”（martingale）概念也获得了独立于随机微积分的表述——第10章正是靠这套哈密顿语言才精确解出了非线性远期利率的鞅测度，并为设计和对冲新的路径依赖工具提供了框架。

第二层：远期利率必须用量子场论。在完全一般的理论里，同一时刻每个期限 $x$ 上的远期利率 $f(t,x)$ 都是独立随机变量，自由度无穷多；把远期利率曲线看作一条随机涨落的曲线，它就是一个二维量子场（quantum field）。场论框架的回报是具体的：高斯场论模型足以刻画远期利率的实证行为；带随机波动率的非线性远期利率可以严格建模；国债的对冲在场论模型下可以有效实施；非线性远期利率的鞅条件与计价物变换（change of numeraire）能用场论技巧精确求解。原书特别强调：全书的实证分析都是在高斯理论中完成的，“刚性”（stiff）场论与市场数据符合得相当好——这为继续研究更复杂的非线性理论提供了动机。

第三层：非线性理论是开放的疆域。无穷多自由度带来非线性场论中潜在的发散（divergence），处理它们的办法是重整化（renormalization）；只有可重整化的理论才是自洽的。原书给后来者留下的核心课题由此清晰：发展解析与数值两方面的计算技术，判定非线性远期利率理论是否可重整化；若可重整化，再发展检验和应用这些理论的实证方法，落实到债务市场。

全书读后建议：把第10--11章合起来看，Baaquie 这本书的方法论主张其实是一句话——金融建模的“合法算符”应当与自由度的数目匹配：单变量问题用布朗运动与二阶微分方程足矣，有限个自由度交给路径积分与哈密顿量，而远期利率这样的场就应当用场论的语言（作用量、传播子、生成泛函、重整化）来提问。读者日后遇到任何利率模型，不妨先问一句“它的自由度是几个、域是什么形状”，再决定该用哪一层的工具——这正是本书比“又一批利率模型”更值得精读的地方。

%TAGS:
